\setcounter{aufgabe}{35}
\einheitenkopf{Einheit 5 von 5 · Diagramme beurteilen und Boxplot}

\begin{aufgabe}{Beurteilen -- Aussagen prüfen. Das Diagramm zeigt, wie viele Bücher eine Bibliothek in fünf Jahren ausgeliehen hat. Kreuze bei a) bis f) an, ob die Aussage stimmt.}
\noindent\saeulen{2020/840,2021/910,2022/880,2023/1050,2024/1200}{1400}{200}{Ausleihen}

\begin{teile}
\teil „2021 wurden 910 Bücher ausgeliehen.“ \hfill \janein
\teil „2022 wurden mehr als 900 Bücher ausgeliehen.“ \hfill \janein
\teil „2024 wurden 1200 Bücher ausgeliehen.“ \hfill \janein
\teil „2020 wurden weniger als 800 Bücher ausgeliehen.“ \hfill \janein
\teil „Von 2023 auf 2024 sind die Ausleihen gestiegen.“ \hfill \janein
\teil „Von 2020 bis 2024 sind die Ausleihen in jedem Jahr gestiegen.“ \hfill \janein
\end{teile}

Rechne bei g) bis j) nach und schreibe deine Rechnung auf.
\begin{teile}
\teil „2024 wurden mehr als doppelt so viele Bücher ausgeliehen wie 2020.“ \\[10pt]
\teil „Von 2020 auf 2023 sind die Ausleihen um ein Viertel gestiegen.“ \\[10pt]
\teil „Jede achte Ausleihe des Jahres 2024 war ein Sachbuch. Das sind $8\,\%$.“ \\[10pt]
\teil In einer Zeitung steht: „Die Ausleihen sind seit 2020 durchgehend gestiegen und haben sich bis 2024 mehr als verdoppelt.“ Prüfe beide Teilaussagen einzeln. \\[14pt]
\end{teile}
\end{aufgabe}

\begin{aufgabe}{Beurteilen -- Achse und Verzerrung. Das Diagramm zeigt die Mitgliederzahl eines Sportvereins.}
\noindent\saeulenab[ymin=400,ymax=460,ystep=10,yfein=5,ylabel=Mitglieder,breite=1.6]{2021/410,2022/420,2023/430,2024/450}

\begin{teile}
\teil Bei welchem Wert beginnt die senkrechte Achse? \leerfeld
\teil Der sichtbare Teil der Säule von 2024 ist fünfmal so hoch wie der von 2021. Um wie viel Prozent ist die Mitgliederzahl von 2021 bis 2024 tatsächlich gestiegen? Runde auf eine Stelle nach dem Komma. \leerfeld[\%]
\teil Erkläre, warum das Diagramm einen falschen Eindruck von der Entwicklung erweckt. \\[12pt]
\teil Skizziere im leeren Gitter dasselbe Diagramm mit einer Achse ab 0 und schreibe in einem Satz auf, was sich am Eindruck ändert.
\end{teile}

\noindent\saeulenab[ymax=500,ystep=100,yfein=50,ylabel=Mitglieder,breite=1.6]{2021/,2022/,2023/,2024/}

\begin{teile}
\teil Jemand sagt: „Wenn es so weitergeht, hat der Verein 2030 mehr als 1000 Mitglieder.“ Beurteile diese Fortschreibung. \\[12pt]
\teil Unter dem oberen Diagramm steht in einer Zeitung die Überschrift „Mitgliederboom: Der Verein wächst rasant“. Prüfe die Aussage, gib die tatsächliche Veränderung an und erkläre, wodurch der Eindruck im Diagramm entsteht. (P10 2019) \\[16pt]
\end{teile}
\end{aufgabe}

\begin{aufgabe}{Beurteilen -- Boxplot lesen. Der Boxplot zeigt die Punktzahlen einer Klasse in einem Test.}
\begin{boxplots}[xmin=0,xmax=20,xstep=2,karo=0.8,xlabel=Punkte]
  \bp{3,7,11,15,19}{Klasse 8a}
\end{boxplots}

\begin{teilezwei}
\tz Minimum \leerfeld & \tz Maximum \leerfeld \\
\tz Median \leerfeld & \tz unteres Quartil \leerfeld \\
\tz oberes Quartil \leerfeld & \tz Spannweite \leerfeld \\
\end{teilezwei}
\begin{teile}
\teil Wie viel Prozent aller Punktzahlen liegen in der Box? \leerfeld[\%]
\teil Wie viel Prozent der Klasse haben mehr als 15 Punkte erreicht? \leerfeld[\%]
\end{teile}
\end{aufgabe}

\begin{aufgabe}{Beurteilen -- Fehler finden.}
In einem Diagramm beginnt die senkrechte Achse bei 200. Die Säule für Mai ist doppelt so hoch wie die für April. Tarek schreibt auf:
\rechnung{\text{Im Mai war der Wert doppelt so groß wie im April.}}
\begin{teile}
\teil Finde den Fehler und benenne ihn. \\[12pt]
\teil Der Wert für April ist 220, der für Mai 240. Um wie viel Prozent ist der Wert gestiegen? Runde auf eine Stelle nach dem Komma. \leerfeld[\%]
\end{teile}
\end{aufgabe}

\begin{aufgabe}{Beurteilen -- Begründen. Rechne nicht.}
\begin{teile}
\teil Begründe, warum die senkrechte Achse bei 0 beginnen sollte, wenn man die Höhen von Säulen miteinander vergleichen will. \\[12pt]
\teil In einem Diagramm fällt eine Linie von 2020 bis 2024 gleichmäßig. Jemand verlängert die Linie bis 2030 und behauptet, dann sei der Wert bei null. Beurteile diese Behauptung. \\[12pt]
\end{teile}
\end{aufgabe}
